\documentclass[11pt]{article}
\usepackage[margin=1.05in]{geometry}
\usepackage{amsmath, amssymb}
\usepackage{booktabs}
\usepackage{multirow}
\usepackage{xcolor}
\usepackage{hyperref}
\hypersetup{colorlinks=true, linkcolor=blue!50!black, urlcolor=blue!50!black}

\newcommand{\code}[1]{\texttt{#1}}
\DeclareMathOperator{\hn}{HN}
\DeclareMathOperator{\IG}{IG}
\DeclareMathOperator{\sd}{sd}
\newcommand{\E}{\mathbb{E}}

\title{\textbf{Failure ledger of the block-1 study}\\[2pt]
\large Which settings and datasets break, through which mechanism,\\
and what each failure does to the Brier score and the CRPS}
\author{Technical note}
\date{31 July 2026}

\begin{document}
\maketitle

\begin{abstract}
This note is the consolidated ledger of every anomaly observed in the
block-1 positive control (settings $4$--$5$, methods i.i.d.\ and AR(2),
datasets $1$--$10$, train $t=1..50$, forecast $t=51..65$). Each affected
cell is classified into one of four pathologies --- sign reflection (R),
positive-branch inflation (P), data-level under-identification (U), and
heavy-tail realization (H) --- with its mechanism stated as a formula and
its measured effect on the two forecast scores. The summary judgment,
which the companion notes quantify in full, is that only unrepaired
reflections with $|\hat\lambda_0|\gtrsim 5$ move the scores materially
(Brier up to $12\times$); every other pathology leaves the paired
score comparison intact.
\end{abstract}

\section{Model, ridge, and scores}

The fitted model, $K=1$ knot, is
\begin{equation}
  y_{st} \;=\; \beta_0 + x_s'\beta + \lambda z_t + \varepsilon_{st}/\sqrt{\tau_t},
  \qquad z_t \ge 0,\; z_t \sim \hn(0, \tau_t^{-1/2}) \text{ with AR(2) copula},
  \label{eq:model}
\end{equation}
with simulation truth $\lambda = 3$, $\beta_0 = 10$, internal
$\tau$-hyperparameters $(a, b) = (3, 8)$, and shared latent AR(2)
coefficients $\phi = (0.80, -0.35)$ in setting~$4$ and $(0.15, 0.80)$ in
setting~$5$ (spectral radii $0.592$ and $0.973$; memory horizons
$h^\ast = 6$ and $109$ days).

The likelihood sees $(\beta_0, \lambda, z)$ only through
$m_t = \beta_0 + \lambda z_t$. For any $(\beta_0^\ast, \lambda^\ast)$ the
reassignment
\begin{equation}
  z_t^\ast = \frac{m_t - \beta_0^\ast}{\lambda^\ast}, \qquad
  z_t^\ast \ge 0 \;\;\forall t,
  \label{eq:ridge}
\end{equation}
reproduces the likelihood exactly. The flat family
\eqref{eq:ridge} has a positive branch ($\lambda^\ast>0$, path shape
kept) and a reflected branch ($\lambda^\ast<0$, path flipped,
$\beta_0^\ast > \max_t m_t$). Identification along the ridge comes only
from the prior on $z$.

Scores are computed per cell on the validation block, means over leads
$1$--$5$ (the pre-registered primary window). CRPS is the sample CRPS
over the posterior predictive draws; the Brier score uses the $q_{95}$
threshold of the cell's own validation block.

\section{The ledger}

Table~\ref{tab:ledger} lists all $13$ cells whose attempt-0 chain failed
the (pre-registered) assertions under the symmetric prior
$\lambda \sim N(0,20)$, together with their fate under the two remedies
that were subsequently tested, and the three datasets whose fitted
$\tau$-shape falls below the $a=2$ moment boundary of
Section~\ref{sec:heavy}. Pathology codes:

\begin{itemize}\itemsep1pt
  \item[\textbf{R}] \textbf{reflection}: chain captured on the
        $\lambda<0$ branch of \eqref{eq:ridge} during burn-in;
  \item[\textbf{P}] \textbf{positive inflation}: same ridge, positive
        branch ($\hat\lambda$ too large, $z$ compressed, sign correct);
  \item[\textbf{U}] \textbf{under-identification}: the $50$-day training
        window realized a nearly flat $z$ path, so $\lambda$ is not
        identified by the data (stable across four seeds and both priors);
  \item[\textbf{H}] \textbf{heavy-tail realization}: fitted $a < 2$, so
        the predictive-spread statistic has no finite sampling variance
        (a property of the dataset, not of any chain).
\end{itemize}

\begin{table}[ht]
\centering
\caption{All affected cells. $\hat\lambda_0$ = attempt-0 posterior mean
under $N(0,20)$; ``guard'' = assertion-guarded reseeding (att = attempts
used, 4 = flagged); ``HN'' = single fit under $\lambda \sim \hn(0,20)$,
same seed, no guard. B/C = truth-recovery and spread assertions,
recorded descriptively. Scores are Brier $q_{95}$, leads 1--5.}
\label{tab:ledger}
\small
\begin{tabular}{llcrclcrr}
\toprule
cell & & type & $\hat\lambda_0$ & guard (att) & HN: $\hat\lambda$, B/C &
\multicolumn{3}{c}{Brier $q_{95}$} \\
\cmidrule(lr){7-9}
 & & & & & & att-0 & guard & HN \\
\midrule
s4 d5  m2 & AR(2)  & R & $-7.28$  & $+3.93$ (2) & $+3.56$, \checkmark\checkmark & .0205 & .0017 & .0016 \\
s4 d5  m1 & i.i.d. & R & $-7.23$  & $+3.52$ (2) & $+3.56$, \checkmark\checkmark & .0265 & .0033 & .0036 \\
s4 d10 m1 & i.i.d. & R & $-4.97$  & $+2.99$ (2) & $+2.84$, \checkmark\checkmark & .0558 & .0226 & .0184 \\
s5 d7  m1 & i.i.d. & R & $-3.99$  & $+2.47$ (2) & $+2.52$, \checkmark\checkmark & .0832 & .0463 & .0465 \\
s4 d4  m2 & AR(2)  & R & $-4.01$  & $+2.19$ (2) & $+2.12$, \checkmark\checkmark & .0010 & .0022 & .0019 \\
s4 d6  m2 & AR(2)  & R+H & $-8.90$ & $+2.86$ (3) & $+3.23$, \checkmark$\times$ & .1378 & .1349 & .1338 \\
s5 d5  m1 & i.i.d. & R & $-4.28$  & $+2.60$ (2) & $+2.58$, \checkmark\checkmark & .0164 & .0166 & .0166 \\
s5 d6  m2 & AR(2)  & P & $+5.40$  & $+3.95$ (2) & $+2.92$, \checkmark\checkmark & .0794 & .0799 & .0808 \\
s5 d8  m1 & i.i.d. & R & $-2.62$  & $+2.29$ (4) & $+2.27$, \checkmark\checkmark & .0000 & .0007 & .0005 \\
s5 d3  m2 & AR(2)  & R & $-15.18$ & $+3.34$ (4) & $+4.97$, $\times$\checkmark & .1088 & .1323 & .1282 \\
\midrule
s5 d2  m1 & i.i.d. & U & $+0.23$  & flag (4)    & $+0.29$, $\times$\checkmark & .0194 & .0194 & .0194 \\
s5 d2  m2 & AR(2)  & U & $+0.55$  & flag (4)    & $+0.57$, $\times$\checkmark & .0190 & .0190 & .0190 \\
s5 d8  m2 & AR(2)  & R & $-2.47$  & flag (4)    & $+3.63$, \checkmark\checkmark & .0000 & .0000 & .0005 \\
\midrule
\multicolumn{9}{l}{\emph{heavy-tail datasets} (property of the data; both
methods, both priors agree on $\hat a$ to two decimals):} \\
s4 d1 & both & H & \multicolumn{6}{l}{$\hat a \approx 1.95$ \quad (truth $a = 3$)} \\
s4 d6 & both & H & \multicolumn{6}{l}{$\hat a \approx 1.51$ \quad ($\to$ the one C exceedance, see \S\ref{sec:heavy})} \\
s5 d5 & both & H & \multicolumn{6}{l}{$\hat a \approx 1.93$} \\
\bottomrule
\end{tabular}
\end{table}

Reflections struck both methods at comparable rates ($6$ of $13$
affected cells are the i.i.d.\ method), which established early that the
pathology belongs to the shared skew-t family of \eqref{eq:model}, not
to the AR(2) time structure.

\section{Mechanisms}

\paragraph{R --- reflection.} Once a chain drifts onto the
$\lambda^\ast < 0$ branch of \eqref{eq:ridge} during burn-in, escape
requires simultaneously moving $\lambda$ across zero, dropping $\beta_0$
by $\approx 20$, and re-inverting the entire $z$ path; every
one-coordinate Gibbs or MH move is individually catastrophic for the
likelihood, so the posterior is bimodal and the chain cannot switch.
Capture is decided by burn-in randomness, which is why reseeding cured
$10$ of $13$ cells.

\paragraph{The $1/|\lambda|$ compression and the forecast law.} On the
ridge $\sd_t(z) = \sd_t(m)/|\lambda|$, so an inflated loading forces an
over-smooth path. Forecasting redraws $z$ from the model law at full
marginal variance, which yields the empirical law
\begin{equation}
  \max_h \sd(\hat y_h) \;\approx\; 1.49\,|\hat\lambda|
  \qquad (R^2 = 0.987 \text{ through the origin}),
  \label{eq:law}
\end{equation}
the single relation that turns a hidden fitting inconsistency into a
visible forecast explosion, and that makes the score damage grow with
$|\hat\lambda_0|$.

\paragraph{The accomplice.} The over-smooth path should be penalized by
the $z$ prior, whose path density per step is
$-\varepsilon_t^2/2\sigma_\varepsilon^2(\phi_z) - \log
\sigma_\varepsilon(\phi_z)$ with innovation scale
$\sigma_\varepsilon(\phi_z) \to 0$ as the AR(2) roots approach the unit
circle. Because $\phi_z$ is estimated, its update reads the compressed
path as evidence of high persistence, drifts toward the unit circle, and
reprices the violation as legal long memory. The dynamics parameter is
the accomplice, not the trigger.

\paragraph{U --- under-identification.} Setting~$5$ has
$h^\ast \approx 109$ days against a $50$-day training window. A window
that realizes a nearly flat $z$ path carries no information about
$\lambda$; dataset~$2$ reproduced $\hat\lambda \approx 0.25$ (i.i.d.)
and $\approx 0.56$ (AR(2)) across four seeds and both priors, stable to
three decimals. No sampler-side or prior-side intervention can create
information the window does not contain.

\paragraph{H --- heavy tails and the $a = 2$ boundary.}
\label{sec:heavy}
With $\sigma^2 = 1/\tau \sim \IG(a, b)$,
\begin{equation}
  \E[\sigma^4] \;=\; \frac{b^2}{(a-1)(a-2)},
  \label{eq:fourth}
\end{equation}
which is infinite for $a \le 2$. The truth $a = 3$ has finite fourth
moment; three datasets nevertheless realized training windows whose
fitted $\hat a < 2$ (Table~\ref{tab:ledger}). Below the boundary the
sample SD of the predictive cloud is an estimator with infinite sampling
variance, so it never stabilizes at any Monte-Carlo size, and a
handful of iterations --- each amplified by the fact that one $\tau_t$
draw scales all $144$ stations at once --- decide whether a spread
threshold is crossed. The one C exceedance at s4 d6 m2 under HN (worst
lead SD $36.9$ against a threshold of $13.1$) is this coin flip, not a
sick chain; the four chains fitted to d6 agree in their posteriors and
in every predictive quantile up to $q_{99.9}$.

\section{Which pathologies touch the scores}

\begin{table}[ht]
\centering
\caption{Score impact by pathology, leads 1--5. ``If shipped'' = the
damage of delivering the attempt-0 chain, measured by bit-for-bit refits
of every replaced chain.}
\label{tab:impact}
\small
\begin{tabular}{llll}
\toprule
type & cells & impact if shipped & impact on paired AR(2)$-$i.i.d. \\
\midrule
R, $|\hat\lambda_0| \ge 5$ & 4 & Brier $8$--$12\times$, CRPS $+40\%$ & material if unprotected \\
R, $|\hat\lambda_0| < 5$   & 6 & small to none (law \eqref{eq:law})   & none measured \\
P                          & 1 & none (.0794 vs .0799)                & none \\
U                          & 2 & none (scores identical)              & none (symmetric across methods) \\
H                          & 3 datasets & none (predictive clouds agree to $q_{99.9}$) & none \\
\bottomrule
\end{tabular}
\end{table}

Three observations close the ledger.

\begin{enumerate}\itemsep2pt
  \item The only score-relevant pathology is an unrepaired large
        reflection. Across the $13$ affected cells, repairing them
        improved the Brier score by $16\%$ and the CRPS by $12\%$ on
        average, with the whole effect concentrated in the
        $|\hat\lambda_0| \ge 5$ rows of Table~\ref{tab:ledger}.
  \item On difficult validation windows an inconsistent fit can
        \emph{outscore} its healthy replacement (s5 d3 m2, twice, under
        both remedies), because overdispersion hedges. Selection must
        therefore never be done on scores.
  \item U and H are properties of the data, are symmetric across the
        compared methods, and cancel in the paired comparison; they
        require reporting, not repair.
\end{enumerate}

Companion notes\, --- \,\code{tex/lambda\_phiz\_ridge} (geometry and law
\eqref{eq:law}), \code{tex/guard\_effect} (score counterfactual of the
guard), and \code{tex/ar2\_remedy\_ledger}, \code{tex/ar2\_score\_findings}
(remedies; main score findings).

\end{document}
